\section{Results (automated pilot)}
\label{sec:results}
All numbers come from run \texttt{pilot-v0.3}: four task families, six conditions and ten candidates per family, with 240 graded pairs. Every pair was graded under harness 0.3.0 at one code revision, with no infrastructure outcomes. Labels are the pre-specified rule's proposals (\S\ref{sec:protocol}), not human labels, so this section is descriptive. Of the 32 held-out candidates, 17 are labeled invalid, 10 valid and 5 ambiguous; the ambiguous ones are excluded from the rates (Table~\ref{tab:labels}). Tables are generated by \texttt{paper/make\_tables.py}; every number is listed in \texttt{paper/evidence\_ledger.md}.

\IfFileExists{generated/results_summary.tex}{\begin{table}[t]\centering\footnotesize
\caption{Verifier reward fidelity on the held-out pool, run \texttt{pilot-v0.3} (4 tasks, 4 families). \textbf{Automated pilot: labels are rule-based proposals (D19), not human labels.} Rates with (num/den) and 95\% task-family bootstrap interval. Reference solution and no-op excluded; ambiguous candidates excluded. $J$ = (1$-$FR)$-$FA.}\label{tab:summary}
\begin{tabular}{lllrrr}\toprule Condition & False acceptance & False rejection & $J$ & Tasks w/ accepted neg. & Infra \\\midrule
C1 initial, plain & 0.35 (6/17) [0.07, 0.71] & 0.20 (2/10) [0.00, 0.67] & +0.45 & 3/4 & 0/32 \\
C2 initial, audited & 0.29 (5/17) [0.00, 0.63] & 0.20 (2/10) [0.00, 0.67] & +0.51 & 2/4 & 0/32 \\
C3 full, plain & 0.00 (0/17) [0.00, 0.00] & 0.00 (0/10) [0.00, 0.00] & +1.00 & 0/4 & 0/32 \\
C4 full, audited & 0.00 (0/17) [0.00, 0.00] & 0.00 (0/10) [0.00, 0.00] & +1.00 & 0/4 & 0/32 \\
B2 full-session LLM & 0.06 (1/17) [0.00, 0.14] & 0.00 (0/10) [0.00, 0.00] & +0.94 & 1/4 & 0/32 \\
B1 repository tests & 0.29 (5/17) [0.11, 0.64] & 0.40 (4/10) [0.00, 0.82] & +0.31 & 3/4 & 0/32 \\
\bottomrule\end{tabular}\end{table}
}{\unfilled{results table}}
\IfFileExists{generated/negatives_table.tex}{\begin{table}[t]\centering\footnotesize
\caption{Accepted invalid candidates by origin and by negative type (accepted/total) (automated pilot). Correction-violating: violates a reference requirement introduced after the first user request.}\label{tab:negatives}
\begin{tabular}{lrrrr}\toprule Condition & Pre-correction state & Near miss & Correction-violating & Initial-violating \\\midrule
C1 initial, plain & 2/2 & 4/15 & 6/16 & 0/1 \\
C2 initial, audited & 2/2 & 3/15 & 5/16 & 0/1 \\
C3 full, plain & 0/2 & 0/15 & 0/16 & 0/1 \\
C4 full, audited & 0/2 & 0/15 & 0/16 & 0/1 \\
B2 full-session LLM & 0/2 & 1/15 & 1/16 & 0/1 \\
B1 repository tests & 0/2 & 5/15 & 5/16 & 0/1 \\
\bottomrule\end{tabular}\end{table}
}{}

\paragraph{Initial-only verifiers reward what the correction ruled out.} Verifiers generated from the initial request (C1) behave as follows (Tables~\ref{tab:summary} and~\ref{tab:negatives}):
\begin{itemize}
\item They accept 6 of the 17 invalid held-out candidates, and every accepted one violates a requirement introduced by the correction.
\item They accept both pre-correction states that the rule labels invalid.
\item They reject 2 of the 10 valid candidates.
\end{itemize}
The same generation procedure given the correction-aware contract (C3) accepts none of the 17 and rejects none of the 10. The paired difference is $+0.35$ in false acceptance and $+0.20$ in false rejection. In the discordant pairs, C1 accepts and C3 rejects six candidates; C1 rejects and C3 accepts two (Table~\ref{tab:comparisons}). The pattern is the same with the audit applied (C2 vs.\ C4). The intervals over four families are wide, and we read them as variability, not as evidence of an effect.

\paragraph{The effect is concentrated where the correction carries information the first request lacks.} In the two families where the correction adds requirements that cannot be read off the first request (\texttt{humanize}, \texttt{more-itertools}), the initial-only verifiers err as follows:
\begin{itemize}
\item In \texttt{more-itertools} they accept the pre-correction state and three of the four invalid near misses.
\item In \texttt{humanize} they accept the pre-correction state and reject both valid alternatives. They also reject the reference solution itself (Table~\ref{tab:sanity}), because the recovered initial-only contract demands the default-output change that the correction retracted.
\end{itemize}
In the other two families the initial-only verifiers err little, for different reasons:
\begin{itemize}
\item In \texttt{tinydb} they do not err at all. The proposer had already inferred, from the first request alone, every requirement the correction later stated (Table~\ref{tab:coverage}).
\item In \texttt{inflection} the plain verifier accepts one invalid near miss (all-uppercase input with hyphens or digit--letter boundaries), which the audited one rejects. The agent's own pre-correction state already satisfied the correction, so there the anchor candidate is valid.
\end{itemize}

\paragraph{The audit changed one outcome.} The predefined audit changed exactly one of the 64 held-out outcomes of the generated-test conditions. On the initial-only contract it caught one near miss (C1 vs.\ C2, $+0.06$ in false acceptance); on the correction-aware contract it changed nothing. Audit:v1 cannot repair a verifier whose contract is wrong. In \texttt{humanize} the audit flagged that the reference fails the initial-only tests, but a contract-faithful revision cannot fix that (decision log D27).

\paragraph{Baselines.} The full-session LLM baseline (B2) accepts 1 of 17 invalid candidates: a near miss that changes rounding of the default output for some inputs. It rejects no valid one. Repository tests at the reference state (B1) are the least faithful condition on both counts. They accept 5 of 17 invalid candidates, because they lack the correction's invariants. They reject 4 of 10 valid ones, because they test reference-specific behaviour the contract does not require, such as temporary-file clean-up on reopen and private hooks.

\paragraph{The headline numbers hide a selection effect.} The automated labels exclude the five near misses whose construction claim the reference checks could not confirm. These are exactly the subtle cases:
\begin{itemize}
\item a stale file only after \texttt{os.chdir} with a relative path;
\item float rounding in the default output;
\item acronyms joined by underscores.
\end{itemize}
In the exploratory analysis that labels constructed candidates by construction intent (Table~\ref{tab:sensitivity}; not pre-specified), the correction-aware verifiers accept 3 of 22 invalid candidates and the full-session LLM baseline 5 of 22. The zero false-acceptance rate of C3 and C4 is therefore conditional on violations that independent checks can detect.
\IfFileExists{generated/results_comparisons.tex}{\begin{table}[t]\centering\footnotesize
\caption{Pre-specified paired comparisons on the same candidates, run \texttt{pilot-v0.2} (automated pilot). Last column: A accepts \& B rejects / A rejects \& B accepts.}\label{tab:comparisons}
\begin{tabular}{lllllrr}\toprule A & B & Metric & A$-$B & 95\% CI & n & discordant \\\midrule
C1 initial, plain & C3 full, plain & FA & +0.71 & [+0.50, +0.80] & 16 & 5/0 \\
C1 initial, plain & C3 full, plain & FR & +0.00 & [+0.00, +0.00] & 16 & 0/0 \\
C2 initial, audited & C4 full, audited & FA & +0.57 & [+0.00, +0.80] & 16 & 4/0 \\
C2 initial, audited & C4 full, audited & FR & +0.00 & [+0.00, +0.00] & 16 & 0/0 \\
C1 initial, plain & C2 initial, audited & FA & +0.14 & [+0.00, +0.50] & 16 & 1/0 \\
C1 initial, plain & C2 initial, audited & FR & +0.00 & [+0.00, +0.00] & 16 & 0/0 \\
C3 full, plain & C4 full, audited & FA & +0.00 & [+0.00, +0.00] & 16 & 0/0 \\
C3 full, plain & C4 full, audited & FR & +0.00 & [+0.00, +0.00] & 16 & 0/0 \\
B2 full-session LLM & C4 full, audited & FA & +0.00 & [+0.00, +0.00] & 16 & 0/0 \\
B2 full-session LLM & C4 full, audited & FR & +0.00 & [+0.00, +0.00] & 16 & 0/0 \\
B2 full-session LLM & C3 full, plain & FA & +0.00 & [+0.00, +0.00] & 16 & 0/0 \\
B2 full-session LLM & C3 full, plain & FR & +0.00 & [+0.00, +0.00] & 16 & 0/0 \\
B1 repository tests & C4 full, audited & FA & +0.43 & [+0.20, +1.00] & 16 & 3/0 \\
B1 repository tests & C4 full, audited & FR & +0.40 & [+0.33, +0.50] & 16 & 0/2 \\
\bottomrule\end{tabular}\end{table}
}{}
\IfFileExists{generated/sanity_table.tex}{\begin{table}[t]\centering\footnotesize
\caption{Sanity checks (excluded from the rates): reference solution and no-op passes per condition.}\label{tab:sanity}
\begin{tabular}{lrr}\toprule Condition & Reference passes & No-op passes \\\midrule
C1 initial, plain & 3/4 & 0/4 \\
C2 initial, audited & 3/4 & 0/4 \\
C3 full, plain & 4/4 & 0/4 \\
C4 full, audited & 4/4 & 0/4 \\
B2 full-session LLM & 4/4 & 0/4 \\
B1 repository tests & 4/4 & 0/4 \\
\bottomrule\end{tabular}\end{table}
}{}
\IfFileExists{generated/sensitivity_table.tex}{\begin{table}[t]\centering\footnotesize
\caption{\textbf{Exploratory, not pre-specified:} rates when constructed candidates are labeled by construction intent (including those the pre-specified rule marked ambiguous).}\label{tab:sensitivity}
\begin{tabular}{lll}\toprule Condition & False acceptance & False rejection \\\midrule
C4 full, audited & 0.18 (2/11) [0.17, 0.20] & 0.00 (0/5) [0.00, 0.00] \\
C3 full, plain & 0.18 (2/11) [0.17, 0.20] & 0.00 (0/5) [0.00, 0.00] \\
C2 initial, audited & 0.64 (7/11) [0.40, 0.83] & 0.00 (0/5) [0.00, 0.00] \\
C1 initial, plain & 0.73 (8/11) [0.60, 0.83] & 0.00 (0/5) [0.00, 0.00] \\
B2 full-session LLM & 0.27 (3/11) [0.17, 0.40] & 0.00 (0/5) [0.00, 0.00] \\
B1 repository tests & 0.55 (6/11) [0.33, 0.80] & 0.40 (2/5) [0.33, 0.50] \\
\bottomrule\end{tabular}\end{table}
}{}

\paragraph{Label uncertainty.} The independent construction checks dispute the proposals of two candidates the rule labels valid: one \texttt{inflection} alternative and the \texttt{inflection} pre-correction state. Both diverge from the reference only on exotic Unicode input, such as modifier apostrophes and the Turkish dotted~İ, and match it on every ASCII probe. The checks also dispute a \texttt{tinydb} alternative, which fails read-after-write only when reopening the file fails. Together with the five ambiguous candidates, 8 of the 32 held-out candidates are queued for human review (Table~\ref{tab:labels}).
\IfFileExists{generated/labels_table.tex}{\begin{table}[t]\centering\footnotesize
\caption{Proposed labels of held-out candidates under rule D19 and the queue for human review (ambiguous, unlabeled, or disagreeing with the independent construction check).}\label{tab:labels}
\begin{tabular}{lrrrrr}\toprule Task & Valid & Invalid & Ambiguous & Unlabeled & Queued \\\midrule
CS-03 & 3 & 2 & 3 & 0 & 3 \\
CS-04 & 2 & 5 & 1 & 0 & 1 \\
\bottomrule\end{tabular}\end{table}
}{}

\paragraph{Contract recovery.} Recovered from the whole session, the full-session contracts cover every active requirement of the reference contract in all four families. The initial-only contracts cover between half and all of them (Table~\ref{tab:coverage}). In \texttt{humanize}, four of the five scored initial-only requirements contradict the reference. A second, independent model mapper agrees closely with the first ($\kappa = 0.92$--$1.00$), but both are the same model.
\IfFileExists{generated/coverage_table.tex}{\begin{table}[t]\centering\footnotesize
\caption{Recovered contracts against the reference contract (model mappings, proposals). Coverage: share of active scored reference requirements that a recovered requirement maps to without contradicting it.}\label{tab:coverage}
\begin{tabular}{llrrrrr}\toprule Task & Contract & Coverage & Unsupported & Contradicting & Coverage (mapper 2) & $\kappa$ \\\midrule
CS-03 & full-session & 1.00 (4/4) & 1/5 & 0 & -- & -- \\
CS-03 & initial-only & 0.75 (3/4) & 1/5 & 0 & -- & -- \\
CS-04 & full-session & 1.00 (6/6) & 1/7 & 0 & -- & -- \\
CS-04 & initial-only & 0.50 (3/6) & 0/3 & 0 & -- & -- \\
\bottomrule\end{tabular}\end{table}
}{}

\paragraph{Harness reliability.} A superseded run under harness 0.1.0 graded 108 of the same pairs, and 107 of them agree with the final run on reward. The one discordant pair is the repository-test grading that the shared-UID defect turned into a ten-minute stall (Table~\ref{tab:reliability}).
\IfFileExists{generated/reliability_table.tex}{\begin{table}[t]\centering\footnotesize
\caption{Agreement of superseded runs with the final run \texttt{pilot-v0.3} on pairs graded in both (pass vs non-pass; exact outcome). Superseded runs used earlier harness versions (D29--D31).}\label{tab:reliability}
\begin{tabular}{lrrr}\toprule Superseded run & Pairs & Reward agreement & Outcome agreement \\\midrule
\texttt{pilot-v0.2} & 108 & 108/108 & 108/108 \\
\texttt{pilot-v0.2-h2} & 4 & 4/4 & 4/4 \\
\bottomrule\end{tabular}\end{table}
}{}
